\section{Literature review}\label{sec:pre_confirmation:literature_review}
Effective abstraction in programming systems requires more than improved libraries and helpers; it requires information hiding, semantic preservation across compilation layers and a reduction in the cognitive burden placed on the programmer~\cite{parnas1972criteria, DiMatteo2024AbstractionHierarchy, yuan2025foundationalabstractions}.
Two complementary forms of abstraction are relevant to evaluating quantum programming systems.
\textit{Functional abstraction} describes hierarchical layers with distinct implementation roles; it governs how a compilation stack is organised and how complexity is distributed across layers.
\textit{Conceptual abstraction} describes the breadth and generality of the concepts a programmer must command to use a layer productively; it governs whether the programmer's mental model is aligned with the problem domain or with the machine.
Classical programming systems have achieved both.
The von Neumann model provided a conceptual foundation whose mental model aligned with human problem-solving idioms, enabling progressive language development from assembly through structured, object-oriented and functional paradigms without requiring programmers to descend the compilation stack~\cite{abelson1996sicp, nunez_corrales2025betterabstractmachines}.
Quantum computing has advanced at the functional level; the conceptual level remains largely unaddressed.
An abstraction hierarchy for quantum computing following the same \textbf{Programming}, \textbf{Execution} and \textbf{Hardware} layers as classical systems has been formally established~\cite{DiMatteo2024AbstractionHierarchy}.
Systematic analysis demonstrates consistent bleed-through between these layers: programming-layer systems require execution-layer and hardware-layer vocabulary as mandatory programmer-facing constructs~\cite{DiMatteo2024AbstractionHierarchy}.
Free value copying, stateful inspection and conditional branching are the classical abstraction contracts that underpin productive programming.
Each is physically constrained in the quantum domain by the no-cloning theorem and related quantum mechanical constraints~\cite{yuan2025foundationalabstractions}.

The quantum programming system landscape includes multiple types of systems and languages offering varying abstraction levels.
Assembly languages such as OpenQASM~\cite{cross2022openqasm3} provide portable, hardware-agnostic circuit specifications; their programmer-facing vocabulary is the execution model by design.
Intermediate representations such as QIR and Q-MLIR~\cite{qiralliance_spec, McCaskey2021MLIRDialectQASM} provide machine-independent substrates for compiler analysis and optimisation; they target toolchain authors rather than application programmers.
Their abstraction boundary does not propagate upward to the programmer's experience.
Frameworks such as Qiskit~\cite{aleksandrowicz2019qiskit} and PennyLane~\cite{bergholm2018pennylane} meet programmers in familiar classical environments, reducing ergonomic barriers without changing the required mental model.
High-level languages including Silq~\cite{bichsel2020silq}, Q\#~\cite{MicrosoftQSharpOverview}, Qmod~\cite{vax2025qmodexpressivehighlevelquantum}, Quff~\cite{wright2024quff} and Qutes~\cite{faro2025qutes} represent the furthest progress toward conceptual abstraction in the current landscape.
Each makes a principled contribution: Silq introduces compiler-managed uncomputation inferred from type annotations; Quff provides runtime function inversion; Qutes introduces scope-bounded uncomputation semantics.
These are advances in functional abstraction at specific, principled points.
The programmer-facing surface model remains quantum-explicit in each case.

The following assessment evaluates representative advanced quantum programming systems against four criteria governing the programming layer~\cite{nunez_corrales2025betterabstractmachines}.
Each records at least one violation.

\paragraph{Symbolic denotational syntax} no circuit elements required to reason about program effects on the state.
\paragraph{Representation-independent data types} no reference to circuit elements in data types.
\paragraph{Classical-quantum regularity} classical and quantum instructions coexist at the same level.
\paragraph{Intrinsic ensemble semantics} shot count elided; distribution types are well-defined at the language level.

\begin{longtable}{p{1.8cm} p{3.3cm} p{3.3cm} p{3.3cm} p{3.3cm}}
    \toprule
    \textbf{System}
    & \textbf{Symbolic denotational syntax.}
    & \textbf{Representation-independent data types.}
    & \textbf{Classical-quantum regularity.}
    & \textbf{Intrinsic ensemble semantics.} \\
    \midrule
    \endhead

    \textbf{Silq}~\cite{bichsel2020silq}
    & \texttt{qfree}, \texttt{const}, \texttt{lifted} and \texttt{mfree} annotations encode circuit-execution semantics as mandatory surface syntax. Both require circuit-model knowledge to apply correctly.
    & \texttt{uint[n]} hard-codes qubit count \texttt{n} into every numeric type. Every declaration couples data representation to a physical resource quantity.
    & The \texttt{!} annotation mandates a programmer-visible stratification between classical and quantum values, with different rules for control flow and function application.
    & Explicit \texttt{measure} is required to obtain classical output. Shot management is external to the language. \\ \addlinespace

    \textbf{Quff}~\cite{wright2024quff}
    & Runtime function inversion (\texttt{invertF}) is a circuit-level surface concept. Phase notation and explicit qubit operations remain at the programming surface.
    & Explicit bit-width specification is required for all quantum types. The programmer must declare and manage qubit count for each variable.
    & Distinct quantum and classical type families are required. Programmers must reason about which operations are valid in each context.
    & Measurement must be explicitly inserted by the programmer. Shot count for statistical reliability is managed outside the language. \\ \addlinespace

    \textbf{Qutes}~\cite{faro2025qutes}
    & \texttt{ref} advances toward quantum-agnostic pass-by-reference. Explicit gate instructions (\texttt{hadamard}, \texttt{mcp}, \texttt{swap}) are still required for algorithms outside the built-in set.
    & Quantum register types carry explicit size or qubit-level structure for operations outside the built-in set.
    & Classical and quantum contexts remain syntactically distinguished. The classical/quantum boundary is programmer-managed rather than compiler-inferred.
    & Scope-bounded uncomputation handles one class of quantum-to-classical transition. General measurement and shot management remain explicit. \\ \addlinespace

    \textbf{Tower}~\cite{yuan2025foundationalabstractions}
    & Data structure operations (push, pop, insert) require no circuit-element reasoning. \texttt{assign}/\texttt{unassign} sequencing and \texttt{with-do} constructs are mandatory surface syntax for any stateful operation, encoding reversibility obligations absent from classical programming.
    & Data structure types (list, stack, queue, set, string) hide internal pointer layout and memory representation from the programmer. Integer widths are bounded by hardware word size, a constraint with classical analogues.
    & The system operates entirely in superposition with no programmer-visible classical/quantum type stratification. \texttt{assign}/\texttt{unassign} sequencing is a quantum-specific surface obligation with no classical analogue.
    & Measurement and shot management are external to the language. Distribution types are not language-level concepts. \\
    \bottomrule
\end{longtable}
\vspace{2em}

A survey of 251 active quantum programmers found that 83\% hold at least five years of classical programming experience and identified limited conceptual quantum knowledge, alongside inconsistent and difficult syntax, among the principal reported barriers to productive use~\cite{ferreira2024qplusage}.
Violations are consistent across advanced quantum programming systems, including the system that advances furthest toward data structure abstraction operating in superposition.
This indicates the ceiling is a structural property of the existing design space rather than a limitation of individual systems.

Two failure modes account for this pattern structurally~\cite{yuan2025foundationalabstractions}.
\textit{Representation exposure} manifests wherever an interface exposes qubit types or requires annotations such as \texttt{const}, \texttt{qfree} or \texttt{Adjoint}.
Quantum state cannot be copied or inspected without disturbing the computation, so no module boundary at which quantum and classical code meet can be made quantum-opaque.
\textit{Semantic mismatch} is visible in the type restrictions documented across the high-level family.
Annotation requirements that make safe uncomputation, inversion and controlled-operation synthesis possible are surface evidence of where classically valid operations have no general quantum equivalent.
The No-Embedding Theorem establishes that classical control flow cannot be realised in superposition because it cannot be made injective~\cite{yuan2025foundationalabstractions}.
Classical abstractions cannot be lifted into superposition and form a complete quantum language; new paradigms and new protocols are required.
The field's own abstract machine analysis confirms that no quantum abstract machine candidate satisfies the criteria for symbolic reasoning about operations without reference to circuit primitives.
No production language is built on a substrate that meets those conditions~\cite{nunez_corrales2025betterabstractmachines}.

No existing quantum abstract machine satisfies the criteria required for a fully separated programming model~\cite{nunez_corrales2025betterabstractmachines}.
The dominant cause lies in the absence of separation between instruction symbols and circuit primitives~\cite{nunez_corrales2025betterabstractmachines}.
This limitation carries upward to the languages built on these machines; current high-level languages remain closer to programming a gate array than to describing intent through shared concepts~\cite{DiMatteo2024AbstractionHierarchy}.
Achieving separation requires an abstraction model that eliminates the leakage of execution-layer and hardware-layer concerns into the programming layer~\cite{DiMatteo2024AbstractionHierarchy}.
Its surface model must be representation-independent~\cite{nunez_corrales2025betterabstractmachines}.
Its conceptual foundations must be built around quantum-mechanical constraints rather than classical abstractions retrofitted to accommodate them~\cite{yuan2025foundationalabstractions}.
This is not merely an engineering shortfall for at least one core abstraction.
The No-Embedding Theorem establishes that lifting classical conditional jumps into superposition is structurally rather than contingently impossible~\cite{yuan2025foundationalabstractions}.
Classical abstraction models, and extensions of them to the quantum setting, therefore cannot close this gap in general.

A language that closes this gap must satisfy eight requirements, adapting the abstract-machine criteria to the language level~\cite{nunez_corrales2025betterabstractmachines}:
\begin{enumerate}
    \item Programs must be expressible without circuit elements at the programming surface.
    \item Data types must remain independent of their physical quantum representation.
    \item Classical and quantum instructions must occupy the same conceptual level, without boundaries the programmer must manage.
    \item Distribution semantics and shot management must be intrinsic to the language.
    \item Module boundaries must not expose quantum-mechanical constraints to callers.
    \item The language must not assume that classical operations remain valid in quantum contexts without constraint, since classical abstractions do not lift unchanged into the quantum domain~\cite{yuan2025foundationalabstractions}.
    \item Semantic intent must be preserved across all compilation layers.
    \item The mental model the programmer must hold must be expressible in problem-space terms, without quantum-mechanical vocabulary~\cite{DiMatteo2024AbstractionHierarchy}.
\end{enumerate}